\documentclass[11pt]{article}
\usepackage[margin=27mm]{geometry}
\usepackage{amsmath,amssymb,amsthm,booktabs,longtable,hyperref}
\newtheorem{lemma}{Lemma}
\title{Exact arithmetic for the Suzuki classification table}
\author{Human Compression Project}
\date{}
\begin{document}
\maketitle
\begin{abstract}
The previously proved classification formulas reduce each numerical row to
three integer factorizations. We acquire certificates for all $180$ integers
occurring at the intended exponents $1\le e\le60$, and substitute them in the
accepted formulas. The numerical conclusions retain the original exact
values and inequalities separately from their corrected values.
\end{abstract}

\section{The arithmetic certificate}
For $q=2^{2e+1}$ the three integers are
\[
 q-1,\qquad q+2^{e+1}+1,\qquad q-2^{e+1}+1.
\]
The classification formulas use their numbers of distinct prime divisors.
Factoring software supplies candidate factorizations; acceptance uses only
integer multiplication, modular exponentiation and greatest common divisors.

\begin{lemma}
Let $N>2$, and suppose the complete prime factorization of $N-1$ has been
certified. If an integer $a$ satisfies
\[
 a^{N-1}\equiv1\pmod N,\qquad
 \gcd\bigl(a^{(N-1)/\ell}-1,N\bigr)=1
 \quad\text{for every prime }\ell\mid N-1,
\]
then $N$ is prime.
\end{lemma}
\begin{proof}
Choose a prime divisor $r$ of $N$. The first congruence makes $a$ invertible
modulo $r$ and makes its order divide $N-1$. For each prime $\ell\mid N-1$,
the greatest-common-divisor condition shows that this order does not divide
$(N-1)/\ell$. Its $\ell$-part is therefore the entire $\ell$-part of $N-1$.
Hence the order is $N-1$. Since it also divides $r-1$, we have $r\ge N$.
As $r\mid N$, equality holds and $N$ is prime.
\end{proof}

This is the classical elementary criterion based on the factorization of
$N-1$; no new primality theorem is claimed. Starting at $2$, recursive
certificates establish every prime used in the three factorizations.
Multiplying the certified prime powers then proves equality with the target
integer. Counting the distinct certified primes gives the required input
to the classification formulas. No probable-prime answer is an acceptance
premise. The PARI/GP factorization routine is used solely to propose missing
factors; the integer reader is independent of its factoring algorithms.

\section{The complete substitution}
The final certificate joins the $134$ previously accepted target
factorizations with the $46$ remaining targets. Their domains are disjoint
and their union is exactly the three displayed integers for each intended
exponent $e=1,\ldots,60$. The completed collection uses $694$ certified primes:
$417$ retained certificates and $277$ added certificates. It reproduces all
$34$ previously completed numerical rows before extending the table to $60$.

Substitution into the already proved classification formulas gives the
following dispositions for the source's total column:
\[
\begin{array}{lr}
\text{printed exact totals recovered}&18\\
\text{printed lower bounds confirmed}&11\\
\text{printed exact totals corrected}&30\\
\text{printed lower bounds refuted}&1.
\end{array}
\]
At $e=39$ the printed lower bound $10^{23}$ is false. The exact total is
\[
                 19128573095168183144089<10^{23}.
\]
There are also two errors in the printed distinct-prime counts: the minus
torus has three prime divisors at $e=43$, rather than two, and four at the
intended $e=60$, rather than three. The last right-hand row is literally
printed with exponent $30$; its position and parameters correspond to
$60$. Both the printed exponent and this contextual correction are retained.

The independent integer acquisition and the complete row data are part of
the accompanying certificate package. The comparison treats every printed
lower bound as an inequality. It neither rounds a corrected exact total to
force agreement nor replaces a printed glyph silently. The example totals
mentioned in Section~9 at $e=3,6,9$ are respectively
\[
2411,\qquad6453727,\qquad18084183211,
\]
so those particular examples and their stated bounds are recovered. The
broader conclusions concerning proper subfields require the corrected
central-cover classification and are addressed there.

\begin{thebibliography}{9}
\bibitem{source} D. Leemans, \emph{The rank 2 geometries of the simple Suzuki
groups $\operatorname{Sz}(q)$}, Beitr\"age zur Algebra und Geometrie
\textbf{39} (1998), 97--120, especially Table~12 and Section~9.
\end{thebibliography}
\end{document}
